\section{Decimosexta iteración de la Fase 5: rangos de la UI y mensajes}
\label{sec:F5I16}

Esta iteración resuelve algunas validaciones pendientes: los campos numéricos de la UI admiten hoy valores mayores de lo que razonablemente admite el juego, y es la verificación quien bloquea cuando ya se ha escrito. Los \texttt{max} de la UI pasan a límites prácticos por campo, separados de lo que admite el formato del GWC; además, los límites se asientan al cargar y al guardar cualquier fichero del proyecto, y la zona de mensajes se revisa para que no conserve errores de una comprobación anterior.

\subsection{Criterios de trabajo}
\label{sec:F5I16Criterios}

\begin{enumerate}
	\item \textbf{El rango manda dos veces.} Cada campo tiene dos límites (ver la tabla   \ref{sec:F5I16Manejo}): el \texttt{Max.GWI} práctico del editor (999 en milisegundos y rangos amplios, 99 en cuentas de ticks y lanzamientos, 255 en el valor inicial de los contadores) y el \texttt{Max.GWC} del formato (255 en los campos de 1 byte; 65534 en los de 2 bytes, porque \texttt{0xFFFF} significa «sin valor»). Siempre \texttt{Max.GWI} $\leq$ \texttt{Max.GWC}: la UI recorta al práctico y la verificación conserva los del formato como red de seguridad.

	\item \textbf{Reproducir antes de corregir.} Prueba que falle con el código actual (la entrada admite el valor fuera de rango) y pase con el corregido. El punto de los mensajes solo se toca si la prueba reproduce la retención; si no, se descarta y se documenta, como marca el criterio del F5-I15.

	\item \textbf{Ningún recorte silencioso.} Toda limitación que cambie un valor en memoria se anuncia: cuando se recorta un valor se informa de modo que el aviso llegue a la pantalla (zona de mensajes de la ficha y barra superior \texttt{\#estado}). El autor siempre sabe que su fichero va a cambiar al guardar.

	\item \textbf{Sin cambios de formato.} No se toca el formato de los cartuchos GWB, GWI ni el GWC binario: solo los límites en la UI, el recorte de los manejadores y el asentamiento de límites al cargar y al guardar. Los cartuchos existentes siguen valiendo.

	\item \textbf{Un cambio, una prueba, una nota.} Cada punto se anota con su prueba.
\end{enumerate}

\subsection{Manejo de los límites}
\label{sec:F5I16Manejo}

Para poder manejar estos límites se implementan estos procesos:

\begin{itemize}
	\item \textbf{Al cargar} cualquier fichero del proyecto, y \textbf{al guardar} cualquiera de ellos, se limitan los valores a los máximos establecidos de sus campos. Al cargar, el recorte deja un aviso en pantalla con lo limitado (anterior criterio 3); al guardar un GWC se limita a los máximos del GWI y, además, a los máximos del tamaño del campo (1 o 2 bytes).
	
	\item \textbf{Redundancia deliberada.} Los dos filtros anteriores se solapan a propósito como decisión de diseño: al cargar ya se recorta a los máximos del GWI, así que el filtro del guardado no debería activarse nunca; se conserva para que un cambio futuro, una ampliación o una mejora del sistema que no contemple los límites no pueda saltárselos. Se documenta aquí para que no se «arregle» esa duplicidad.
	
	\item \textbf{Nueva función \texttt{revisarValores}} que limita los valores en memoria. Admite un parámetro \texttt{tipo}, uno de los tipos de fichero del proyecto, de modo que un GWS se limita con sus máximos, un GWI con los suyos prácticos y un GWC con los del tamaño de sus campos. Para valorar cada campo llamará a la función auxiliar \texttt{limitarValor} con tres parámetros:
		\begin{itemize}
		\item \texttt{valor} a limitar (devuelve el valor ajustado).
		\item \texttt{maximo} que admite el campo
		\item \texttt{minimo} que admite (opcional; si no se indica, es cero).
		\end{itemize}
	
	\item \textbf{\texttt{revisarValores}} se llama después de cargar cualquier fichero y siempre antes de guardar cualquier fichero, de manera que los valores en memoria sean coherentes en pantalla tras cargar y correctos siempre al generar. Vale para los tres tipos: el GWS no puede manejar un puerto 65000 para un botón, una pantalla de 10$\times$15 o teclas fuera del rango \texttt{A..Z} \texttt{0..9}.
	\item \textbf{El nulo no es cero.} Los campos opcionales del Bloque 1 distinguen «sin valor» (memoria: \texttt{null}; en pantalla: campo vacío, mostrado como \texttt{-1} o \texttt{FFFFh} donde se presente en crudo; en el GWC: \texttt{0xFFFF}) de un valor legal \texttt{0}. \texttt{limitarValor} respeta esa diferencia: no toca un \texttt{null} y no lo convierte jamás en \texttt{0}.
	\item \textbf{Una sola tabla de límites.} Mantenemos una tabla interna única con \texttt{campo} $\rightarrow$ \{ \texttt{min}, \texttt{maxGWI}, \texttt{maxGWC} \}, de forma que \texttt{revisarValores(tipo, valores)} solo elige columna y el proceso sea sencillo de ampliar. 
\end{itemize}

\texttt{Max.GWI} es el límite práctico del editor y \texttt{Max.GWC} el del formato binario. Los campos de lista no admiten recorte (son \texttt{<select>} o validación) y los máximos del nulo (\texttt{0xFFFF} en 2 bytes, \texttt{0xFF} en 1) no se alcanzan con estos valores. 

Hay un límite en el \textbf{tamaño mínimo de la pantalla}, para que sea visible debe ser de $240\times320$; si es menor se dará un aviso pero dejará continuar. El máximo sería de 1024x600, aunque puede llegar a 1920x1080 usando un interfaz MIPI-DSI, pero no es apropiado para una pantalla de un Game \& Watch. La tabla completa es esta:

% --- TABLA 1: GWI / GWC ---
\begin{table}[H]
	\centering
	\begin{tabularx}{\textwidth}{l X l l l}
		\toprule
		\textbf{Campo} & \textbf{Uso} & \textbf{min} & \textbf{\makecell[l]{Max.\\GWI}} & \textbf{\makecell[l]{Max.\\GWC}} \\
		\midrule
		\texttt{valorInicial} & Valor de arranque del contador y órdenes vida/muerte & 0 & 255 & 255 \\
		\texttt{secTicks} & Ticks entre lanzamientos de la Secuencia & 1 & 99 & 255 \\
		\texttt{pausaSim} & Pausa del simulador (ms) & 1 & 999 & 65534 \\
		\texttt{pausaEsp} & Pausa del sistema externo (ms) & 1 & 999 & 65534 \\
		\texttt{accCada} & Ticks entre reducciones de la aceleración & 1 & 99 & 65534 \\
		\texttt{accResta} & Reducción de la aceleración (ms) & 1 & 999 & 65534 \\
		\texttt{accMinimo} & Pausa mínima de la aceleración (ms) & 0 & 999 & 65534 \\
		\texttt{aumentarCada} & Lanzamientos entre ampliaciones de rutas & 1 & 99 & 65534 \\
		\texttt{fin.mecanismo} & Mecanismo de fin (lista) & --- & --- & \makecell[l]{0--3 \\ \texttt{nulo}} \\
		\texttt{fin.contador} & Contador interno del fin (lista) & 1 & 4 & 4 \\
		\texttt{contadorInterno} & Contador interno de un segmento contador & 1 & 4 & 4 \\
		\texttt{resolucion} & Pantalla declarada \{ancho, alto\}\footnotemark & 1 & 999 & 65534 \\
		\texttt{nSprites}, \texttt{nSegmentos} & Conteos del Bloque 0 & 0 & 255 & 255 \\
		\texttt{x}, \texttt{y} de segmento & Posición dentro de la resolución & 0 & \makecell[l]{999 \\ (resol.)} & 65534 \\
		\bottomrule
	\end{tabularx}
	\caption{Límites para los ficheros GWI/GWC}
	\label{tab:limites_gwi_gwc}
\end{table}
\footnotetext{Fuera del tamaño mínimo y máximo de pantalla se dará un aviso pero dejará continuar.}

\begin{table}[H]
	\centering
	\begin{tabularx}{\textwidth}{l X c c}
		\toprule
		\textbf{Campo} & \textbf{Uso} & \textbf{Mín.} & \textbf{Máx.} \\
		\midrule
		\texttt{estilo} & Estilo de máquina (lista 2/4/5) & --- & --- \\
		\texttt{resolucion} & Pantalla del hardware \{ancho, alto\}\footnotemark & 1 & 1024 \\
		\texttt{botones[].numero} & Número de botón (fijo por estilo) & --- & --- \\
		\texttt{botones[].puerto} & Pin del botón en la máquina real & 0 & 99 \\
	    \texttt{botones[].tecla} & Tecla del simulador (\texttt{A..Z} + \texttt{0..9}) & --- & --- \\
		\bottomrule
	\end{tabularx}
	\caption{Límites para los ficheros de sistema GWS.}
	\label{tab:limites_gws}
\end{table}
\footnotetext{Fuera del tamaño mínimo y máximo de pantalla se dará un aviso pero dejará continuar.}


\begin{table}[H]
	\centering
	\begin{tabularx}{\textwidth}{l X l l}
		\toprule
		\textbf{Campo} & \textbf{Uso} & \textbf{Mín} &  \textbf{Máx.} \\
		\midrule
		\texttt{resolucion} & Pantalla opcional del arte\footnotemark & 1 & 1024 \\
		\texttt{ColorFondo} & Color \texttt{\#RRGGBB} (lista/validado) & --- & --- \\
		\bottomrule
	\end{tabularx}
	\caption{Límites para los ficheros de arte GWB.}
	\label{tab:limites_gwb}
\end{table}

\footnotetext{Fuera del tamaño mínimo y máximo de pantalla se dará un aviso pero dejará continuar.}


En la columna \texttt{Max.GWC}: los campos de 2 bytes admiten 65534 porque \texttt{0xFFFF} está reservado como «sin valor»; los de 1 byte admiten 255 con \texttt{0xFF} para los opcionales del Bloque 3. En \texttt{valorInicial} el límite práctico coincide con el del formato. En el GWS, «hardware» anticipa el límite que impondrá el firmware de la Fase 6 (1 byte por puerto); donde hoy no se valida nada, \texttt{revisarValores} pasa a hacerlo.

En los campos de 1 byte el generador ya recorta (\texttt{BaseCartucho.js} escribe \texttt{secTicks} con \texttt{Math.min(0xFF, ...)}). Los \texttt{max} de hoy (999 y 9999) no corresponden a ningún límite del formato ni a una decisión práctica: el navegador permite escribir valores que solo la verificación caza después. Fuera de esta tabla quedan los demás campos de lista, que ya hoy son \texttt{<select>} y no admiten tecleo; no hay \texttt{max} que ajustar.

\subsection{Detalles importantes}
\label{sec:F5I16Detalles}

Hay algunos detalles que considerar al limitar el manejo de estos valores máximos:

\begin{itemize}
	\item El atributo \texttt{max} del HTML no impide escribir a mano un valor mayor (solo gobierna las flechas del campo y la validación de formulario). Hoy los manejadores \texttt{change} recortan por abajo (\texttt{Math.max(minimo, ...)}) pero no por arriba, así que un valor teclado puede quedar en la propiedad sin recortar. Cada punto incluye el recorte superior en el manejador, no solo el atributo.
	\item El recorte superior en el manejador (escribir un valor mayor lo trunca al máximo, con su aviso) es la opción utilizada en el manejo del interfaz.
	\item Un GWI editado a mano con \texttt{pausaSim} de 5000 debe ser válido y generable. El paso a GWC verifica que los límites del formato (255 y 65534) no se sobrepasen como red de seguridad, pero los valores en memoria se limitan a los máximos definidos para el GWI (con su aviso), y lo hace de forma redundante con el guardado: decisión de diseño del apartado anterior.
	\item La limpieza de mensajes cubre la zona de la ficha y también la barra superior \texttt{\#estado}.
\end{itemize}

\subsection{Cambios necesarios en el código}
\label{sec:F5I16Codigo}

\begin{xltabular}{\textwidth}{l c >{\RaggedRight\arraybackslash}X >{\RaggedRight\arraybackslash}X}
	\toprule
	\textbf{Módulo} & \textbf{\#} & \textbf{Qué se hará} & \textbf{Prueba / estado} \\
	\midrule
	\endhead
	BaseCartucho & M1 & Nueva tabla de límites (\texttt{campo} $\rightarrow$ \texttt{min}/\texttt{maxGWI}/\texttt{maxGWC}) y funciones \texttt{revisarValores(tipo, valores)} y \texttt{limitarValor(valor, maximo, minimo)}; \texttt{limitarValor} informa de cada valor recortado & Caja con \texttt{pausaSim} 5000: la memoria queda en 999 y aparece el aviso \\
	BaseCartucho y BaseDirector & M2 & Llamar a \texttt{revisarValores} tras cargar cada fichero del cartucho (GWB/GWI/GWC) y al validar el setup (GWS) & Cargar un GWS a mano con \texttt{puerto} 65000: queda 255 con aviso; cargar un GWI a mano con \texttt{pausaSim} 5000: queda 999 con aviso \\
	BaseCartucho y BaseSetup & M3 & Llamar a \texttt{revisarValores} antes de guardar GWI y GWC (BaseCartucho) y GWS (BaseSetup); el filtro del GWC aplica máximo GWI y máximo de tamaño de campo & Guardar el GWC de un GWI a mano con \texttt{pausaSim} 5000: el GWC lleva 999 y la equivalencia en memoria sigue verde \\
	BaseSetup/BaseDirector & M4 & GWS: \texttt{puerto} limitado a 0--255 (hoy sin validación), \texttt{tecla} ampliada a \texttt{A..Z} + \texttt{0..9} (hoy solo letra en \texttt{validarSetup}) y \texttt{resolucion} con tope 999 & Puerto 65000 queda 255; la tecla «5» se acepta y «AA» se rechaza \\
	Editor & R1 & Campo «Valor inicial» de los contadores: \texttt{max} pasa de 9999 a 255 y el manejador \texttt{change} recorta por arriba & Escribir 256 en el campo deja el valor en 255 \\
	Editor & R2 & Campo «Ticks entre lanzamientos» de la Secuencia: \texttt{max} pasa de 999 a 99 (cuenta ticks), con recorte en el manejador & Escribir 100 deja 99 \\
	Editor & R3 & Campos de 2 bytes del panel del reloj (construidos por \texttt{campoNumero}): \texttt{pausaSim}, \texttt{pausaEsp}, \texttt{accResta} y \texttt{accMinimo} (milisegundos) pasan de 9999 a 999; \texttt{accCada} y \texttt{aumentarCada} (ticks/lanzamientos) pasan de 9999 a 99. Recorte en el manejador en todos & Escribir 1000 deja 999 en las pausas; escribir 100 deja 99 en las cuentas de ticks \\
	Editor & R4 & La zona de mensajes no conserva errores viejos: al aplicar la ficha del cartucho con éxito (y al generar) lo mostrado corresponde al estado actual, en la ficha y en \texttt{\#estado}. \textbf{Punto condicional}: solo se corrige si la prueba reproduce la retención & Prueba: forzar un error, aplicar bien, y comprobar que no queda error anterior visible \\
	\bottomrule
\end{xltabular}

\subsection{La zona de mensajes}
\label{sec:F5I16Mensajes}

Los mensajes viven hoy en la barra superior \texttt{\#estado} (\texttt{BaseDirector.avisar}/\texttt{avisarError}) y en la zona de mensajes de la ficha (\texttt{BaseFicha.\_pintarMensajes}), que se repinta en cada verificación y se limpia al poblar o reiniciar la ficha. El caso a reproducir: un error de una verificación fallida que siga visible después de que una aplicación posterior con éxito lo haya dejado sin sentido. Si la prueba no llega a reproducir retención alguna, el punto se descarta y se documenta. La limpieza cubre la zona de la ficha y también la barra superior, y es también por ahí donde deben salir los avisos de \texttt{limitarValor} (criterio 3).

\subsection{Pruebas}
\label{sec:F5I16Pruebas}

En la suite del editor (\texttt{pruebas/pruebasEditor.js}, Playwright): una caja por bloque de campos, cada una escribe un valor fuera de rango en el campo real del panel y comprueba el valor que queda tras el \texttt{change} (255 en \texttt{valorInicial}; 99 en \texttt{secTicks}, \texttt{accCada} y \texttt{aumentarCada}; 999 en \texttt{pausaSim}, \texttt{pausaEsp}, \texttt{accResta} y \texttt{accMinimo}), y que el aviso de recorte es visible. Para R4, la caja fuerza un error de ficha, aplica con éxito y comprueba ficha y \texttt{\#estado}. Para M2/M3, cargas de ficheros ajustados a mano (GWI con \texttt{pausaSim} 5000, GWS con puerto 65000 y teclas con dígitos): memoria limitada, aviso en pantalla y GWC guardado dentro del formato. En la suite \texttt{pruebasGWI\_GWC} la equivalencia con los ficheros limitados debe seguir en verde. Al cierre se re-lanzan las cuatro suites (\texttt{pruebasMotor}, \texttt{pruebasDirector}, \texttt{pruebasEditor}, \texttt{pruebasGWI\_GWC}); no hace falta red de seguridad de partidas: los límites no tocan el orden del tick.
